\newcommand{\docTopicInc}{
\begin{figure}[h]
\centering
\resizebox{\linewidth}{!}{%
\begin{tikzpicture}[
  font=\small,
  >=Latex,
  node distance=14mm and 16mm,
  box/.style={draw, rounded corners=2pt, thick, minimum height=9mm, align=center},
  arrow/.style={->, thick},
  label/.style={font=\small\bfseries},
]

% -------------------------
% Left: document corpus (cylinder)
% -------------------------
\node (corpus) [cylinder, shape border rotate=90, aspect=1.5,
               draw, thick, minimum height=18mm, minimum width=16mm] {};
\node[below=1mm of corpus] {\textbf{document}};
\node[below=4mm of corpus.south] {\textbf{corpus}};

% -------------------------
% Topic model text + arrow
% -------------------------
\node (tm) [right=24mm of corpus, align=center] {\textbf{Topic model}\\[-1mm]\textbf{LDA}};
\draw[arrow] (corpus.east) -- (tm.west);

% -------------------------
% Stack of topic vectors
% -------------------------
% back plates (to look like a stack)
\node (stack3) [draw, thick, rounded corners=2pt, minimum width=40mm, minimum height=14mm,
                right=18mm of tm, xshift=3mm, yshift=3mm, fill=white] {};
\node (stack2) [draw, thick, rounded corners=2pt, minimum width=40mm, minimum height=14mm,
                right=18mm of tm, xshift=1.5mm, yshift=1.5mm, fill=white] {};
\node (stack1) [box, minimum width=40mm, minimum height=14mm, right=18mm of tm, fill=white] 
  {$\{t_{i0},\ldots,t_{in}\}\in\mathbb{R}^{n}$};

\draw[arrow] (tm.east) -- (stack1.west);

\node[below=1mm of stack1] {\textbf{topic vector}};
\node[below=4mm of stack1.south] {$\forall{\text{ document }i}$};

% -------------------------
% Delta arrow
% -------------------------
\node (delta) [right=18mm of stack1, label] {$\delta$};
\draw[arrow] (stack1.east) -- (delta.west);

% -------------------------
% Right: docs x topics table
% -------------------------
\node (tab) [right=14mm of delta, inner sep=0pt] {%
\begin{tikzpicture}[font=\small, thick]
  % Table parameters
  \def\w{15mm}   % col width
  \def\h{9mm}    % row height
  \def\ncols{4}  % header has: t1 t2 ... tn (we'll draw 4 columns)
  \def\nrows{5}  % docs: d1 d2 ... dl (we'll draw 5 rows total incl header lines)

  % Outer frame (docs column + topic columns)
  \draw (0,0) rectangle ({(1+\ncols)*\w},{\nrows*\h});

  % Vertical lines (separate docs column and topic columns)
  \draw (\w,0) -- (\w,{\nrows*\h});
  \foreach \c in {2,...,\ncols} {
    \draw ({\c*\w},0) -- ({\c*\w},{\nrows*\h});
  }

  % Horizontal lines
  \foreach \r in {1,...,\numexpr\nrows-1\relax} {
    \draw (0,\r*\h) -- ({(1+\ncols)*\w},\r*\h);
    }


  % Slanted top-left "topics/docs" header split
  \draw (0,{\nrows*\h}) -- (\w,{\nrows*\h-\h});
  \node[anchor=west] at (5mm,{\nrows*\h-2mm}) {\textbf{topics}};
  \node[anchor=west] at (1mm,{\nrows*\h-\h+2mm}) {\textbf{docs}};

  % Column headers for topics
  \node at ({1.4*\w},{\nrows*\h-0.5*\h}) {$t_1$};
  \node at ({2.4*\w},{\nrows*\h-0.5*\h}) {$t_2$};
  \node at ({3.4*\w},{\nrows*\h-0.5*\h}) {$\cdots$};
  \node at ({4.4*\w},{\nrows*\h-0.5*\h}) {$t_n$};

  % Row labels (docs)
  \node at ({0.4*\w},{\nrows*\h-1.5*\h}) {$d_1$};
  \node at ({0.4*\w},{\nrows*\h-2.5*\h}) {$d_2$};
  \node at ({0.4*\w},{\nrows*\h-3.5*\h}) {$\vdots$};
  \node at ({0.4*\w},{\nrows*\h-4.5*\h}) {$d_\ell$};

  % Example weights
  \node at ({1.4*\w},{\nrows*\h-1.5*\h}) {$w_{1,1}$};
  \node at ({2.4*\w},{\nrows*\h-1.5*\h}) {$w_{1,2}$};
  \node at ({3.4*\w},{\nrows*\h-1.5*\h}) {$\cdots$};
  \node at ({4.4*\w},{\nrows*\h-1.5*\h}) {$w_{1,n}$};
  \node at ({1.4*\w},{\nrows*\h-2.5*\h}) {$w_{2,1}$};
  \node at ({2.4*\w},{\nrows*\h-2.5*\h}) {$w_{2,2}$};
  \node at ({3.4*\w},{\nrows*\h-2.5*\h}) {$\cdots$};
  \node at ({4.4*\w},{\nrows*\h-2.5*\h}) {$w_{2,n}$};
  \node at ({1.4*\w},{\nrows*\h-3.5*\h}) {$\vdots$};
  \node at ({2.4*\w},{\nrows*\h-3.5*\h}) {$\vdots$};
  \node at ({3.4*\w},{\nrows*\h-3.5*\h}) {$\ddots$};
  \node at ({4.4*\w},{\nrows*\h-3.5*\h}) {$\vdots$};
  \node at ({1.4*\w},{\nrows*\h-4.5*\h}) {$w_{\ell,1}$};
  \node at ({2.4*\w},{\nrows*\h-4.5*\h}) {$w_{\ell,2}$};
  \node at ({3.4*\w},{\nrows*\h-4.5*\h}) {$\cdots$};
  \node at ({4.4*\w},{\nrows*\h-4.5*\h}) {$w_{\ell,n}$};
\end{tikzpicture}
};

\draw[arrow] (delta.east) -- (tab.west);

\end{tikzpicture}
}
\caption{Deriving document-topic incidence matrix from topic model.}
\label{fig:doc_topic_inc}
\end{figure}}

\subsection{Topic Model}
\label{sec:topic_model}
The data-driven topic-model-based approach involves three steps: constructing a topic model, building a document–topic incidence, and filtering concepts via TITANIC~\cite{gutekunst2025conceptual}.

\subsubsection{Topic Model Construction}
A topic model is a generative model for documents, where documents are mixtures of $K$ topics and topics are mixtures of words (i.e., multinomial distribution over a word vocabulary).
Topic models rely on the bag-of-words assumption, i.e., ignoring the order of words.
We employed \texttt{gensim}'s \ac{lda}\footnote{\url{https://radimrehurek.com/gensim/models/ldamodel.html} (30.12.2025)}.
\ac{lda} is a widely used unsupervised statistical bayesian topic model~\cite{alghamdi_survey_2015}.
For a corpus of $N$ documents, \ac{lda} produces a document-topic distribution matrix $\Theta \in \mathbb{R}^{N \times K}$, where $K$ is the number of topics and $\theta_{ij}$ represents the probability that document $i$ is associated with topic $j$.

To construct the formal context, document-topic probabilities are binarized using a threshold $\delta$. 
To select an appropriate threshold, we analyze how different threshold values affect the resulting incidence density. 
Figure~\ref{fig:incidence_density} illustrates this relationship for the BankSearch dataset.

\begin{figure}[t]
    \centering
    \includesvg[width=0.7\textwidth]{results/lda/incidence_density_vs_threshold.svg}
    \caption{\ac{lda} document-topic incidence density for different threshold $\delta$ values for the BankSearch dataset. The elbow criterion suggests an optimal threshold around 0.05, yielding an incidence density slightly above 0.2.}
    \label{fig:incidence_density}
\end{figure}


\subsubsection{Document–Topic Incidence}
Given threshold $\delta$, document $d$ is assigned topic $t$ if and only if $\theta_{dt} \geq \delta$. 
As displayed in Figure~\ref{fig:doc_topic_inc}, this yields a binary incidence relation $I \subseteq D \times T$, where $(d,t) \in I$ if and only if topic $t$ is sufficiently representative of document $d$. 
The threshold $\delta$ controls the sparsity of the resulting formal context; lower values produce denser contexts with more document-topic associations, while higher values yield sparser contexts with only the most dominant topics retained.
The binary incidence relation $I$ defines a formal context $\mathbb{K} = (G, M, I)$, where $G$ is the set of documents (i.e., objects), $M$ is the set of topics (i.e., attributes), and $I$ specifies which topics apply to which documents.       

\docTopicInc{}

\subsubsection{Concept Filtering via TITANIC}

% copied from other paper
Computing complete concept lattices of $\mathbb{K}$ for large datasets presents a significant challenge, as the size of a concept lattice can grow exponentially with respect to the size of the context.
Consequently, complete visualizations of large concept lattices may obscure important hierarchical relationships within the data, hindering human interpretation.
To address this issue, \cite{stumme_titanic_2002} introduced the TITANIC algorithm, which computes iceberg concept lattices.
An iceberg concept lattice includes only the top-level concepts,
specifically those corresponding to frequent attribute sets.

The authors define a minimum support threshold,
$\mathtt{min\_supp} \in [0,1]$, below which attribute sets are not considered frequent.
An attribute set $B \subseteq M$ is deemed frequent in $\mathbb{K}$ if its support $$\operatorname{supp}(B) = \frac{\left| B' \right|}{\left| G \right|}$$ meets or exceeds the \texttt{min\_supp} threshold.
A concept $(A, B)$ is considered frequent if its intent $B$ constitutes a frequent attribute set.
The iceberg concept lattice of a context $\mathbb{K}$ comprises all such frequent concepts.
%  (cf. Figure~\ref{fig:titanic}).
% \displayTITANIC{}

Intuitively, this restriction retains only those document-topic combinations where the topic set is representative of at least $\mathtt{min\_supp} \times 100\%$ of the corpus.
The choice of minimum support significantly influences the granularity of the extracted hierarchy. 
Figure~\ref{fig:concepts_support} demonstrates this relationship for the BankSearch dataset.
Based on this analysis, we select $\texttt{min\_supp} = 0.05$ to balance hierarchy richness with conceptual interpretability, yielding approximately 30 concepts for the BankSearch dataset.
% The resulting iceberg concept lattice is visualized in Figure~\ref{fig:iceberg_concept_lattice}.


\begin{figure}[htbp]
    \centering
    \includesvg[width=0.7\textwidth]{resources/banksearch/topic_model/plots/concepts_vs_support.svg}
    \caption{Number of concepts in the iceberg lattice as a function of minimum support threshold for the BankSearch dataset. The curve exhibits the expected monotonic decrease, with support values below 0.15 yielding fewer than 10 concepts. We select $\mathtt{min\_supp} = 0.05$, resulting in approximately 30 concepts.}
    \label{fig:concepts_support}
\end{figure}

